%----------------------------------------------------------------------------------------
% Shared settings of both CVs (cv.tex and cv_chinese.tex input this file).
% The language is chosen by the wrapper before this file is read:
%   \EN{...} prints its text only in the English CV,
%   \ZH{...} prints its text only in the Chinese CV.
% CV text does not go here: it lives in cv_content.tex.
%----------------------------------------------------------------------------------------

\usepackage[dvipsnames]{xcolor}
\usepackage[left=0.5in,top=0.6in,right=0.6in,bottom=0.6in]{geometry}
\usepackage{xurl}
\usepackage{needspace} % keep subheadings with their first entry
\usepackage[colorlinks=true,urlcolor=Blue,linkcolor=Blue]{hyperref}

%----------------------------------------------------------------------------------------
%	BIBLIOGRAPHY  (publications.bib and talks.bib, in English for both CVs)
%----------------------------------------------------------------------------------------

\usepackage[
  backend=bibtex,   % biber is broken on this machine; plain bibtex is enough here
  style=phys,
  articletitle=true,
  biblabel=brackets,
  eprint=true,
  sorting=none,        % printed in the order entries appear in the .bib files (keep them newest first)
  defernumbers=true,   % number each list separately
  maxnames=1, minnames=1,
  giveninits=true,
]{biblatex}
\addbibresource{publications.bib}
\addbibresource{talks.bib}

% Bold my name wherever it appears
\renewcommand*{\mkbibnamefamily}[1]{\ifdefstring{\namepartfamily}{Zeng}{\textbf{#1}}{#1}}
\renewcommand*{\mkbibnamegiven}[1]{\ifdefstring{\namepartfamily}{Zeng}{\textbf{#1}}{#1}}

% "Contributions:" line under an entry (from the bib `addendum` field, which is in English)
\DeclareFieldFormat{addendum}{\textit{Contributions:} #1}
\renewbibmacro*{addendum+pubstate}{\setunit{\newline}\printfield{addendum}\newunit\printfield{pubstate}}
\renewcommand*{\finentrypunct}{}  % no period at the end of entries
% arXiv numbers never break across lines: the phys style's format, with the number in a \mbox
\makeatletter
\DeclareFieldFormat{eprint:arxiv}{%
  \ifhyperref
    {\href{https://arxiv.org/\abx@arxivpath/#1}{%
        arXiv\addcolon
        \mbox{\nolinkurl{#1}}%
        \iffieldundef{eprintclass}
          {}
          {\addspace\UrlFont{\mkbibbrackets{\thefield{eprintclass}}}}}}
    {arXiv\addcolon
      \mbox{\nolinkurl{#1}}%
      \iffieldundef{eprintclass}
        {}
        {\addspace\UrlFont{\mkbibbrackets{\thefield{eprintclass}}}}}}
\makeatother
% ...and a reference line may stretch a little instead of running into the margin
\AtBeginBibliography{\setlength{\emergencystretch}{1.5em}}
% Talks (@misc): bold title / italic venue, date / link. Author is omitted -- it's always me.
\DeclareFieldFormat[misc]{title}{\textbf{#1}}
\DeclareFieldFormat[misc]{howpublished}{\textit{#1}}
\DeclareBibliographyDriver{misc}{%
  \ifkeyword{invited}{\EN{\textbf{Invited talk:} }\ZH{\textbf{特邀报告：}}}{}%
  \printfield{title}\newline
  \printfield{howpublished}\setunit{\addcomma\space}%
  \usebibmacro{date}%
  \iffieldundef{url}{}{\newline\printfield{url}}%
  \usebibmacro{finentry}}

\setlength{\bibitemsep}{3pt}
\defbibheading{cvsub}[]{\needspace{4\baselineskip}\vspace{2pt}\textbf{\EN{\itshape}#1}\par\vspace{1pt}}
% No visible heading; \leavevmode is needed so the labels of a list placed right
% after the section's \item[] don't pile up
\defbibheading{cvnone}[]{\leavevmode\par\vspace{-\baselineskip}}

%----------------------------------------------------------------------------------------
%	HEADER
%----------------------------------------------------------------------------------------

% TODO: Chinese name for the Chinese CV, e.g. \ZH{曾XX \quad }Jiancheng (J.C.) Zeng
\name{Jiancheng (J.C.) Zeng}
\address{\\ \href{mailto:JCZeng1412@gmail.com}{JCZeng1412@gmail.com} \\ 857-930-3319 \\ \EN{Boston, MA}\ZH{美国马萨诸塞州波士顿}}

%----------------------------------------------------------------------------------------
%	BUILDING BLOCKS FOR cv_content.tex
%----------------------------------------------------------------------------------------

% Blue section titles (upper case in English)
\renewenvironment{rSection}[1]{
  \sectionskip
  \textcolor{Blue}{\bf #1}
  \sectionlineskip
  \hrule
  \begin{list}{}{\setlength{\leftmargin}{1.5em}}
  \item[]
}{
  \end{list}
}

% \begin{cvsection}{English title}{中文标题} ... \end{cvsection}
% (\needspace: a title never sits alone at the bottom of a page)
\newenvironment{cvsection}[2]{\needspace{5\baselineskip}\begin{rSection}{\EN{\MakeUppercase{#1}}\ZH{#2}}}{\end{rSection}}

% Comma between fields: ", " in English, "，" in Chinese
\newcommand{\cvsep}{\EN{, }\ZH{，}}

% \cvposition{job title}{institution}{dates}{small second line}{place}
\newcommand{\cvposition}[5]{{\bf #1}\cvsep #2 \hfill {\em #3} \\ {\small #4} \hfill #5}

% \cvschool{school}{place}{dates}{degree lines, separated by \\}
\newcommand{\cvschool}[4]{{\bf #1}\cvsep #2 \hfill {\em #3} \\ #4}

% Keep a group, project or teaching title on the same page as its first bullet
\BeforeBeginEnvironment{rSubsubsectionTitle}{\needspace{5\baselineskip}}
\BeforeBeginEnvironment{rSubsubsection}{\needspace{3\baselineskip}}
\BeforeBeginEnvironment{rSubsection}{\needspace{3\baselineskip}}

% \cvlabel{English label}{中文标签}: italic label at the start of a skills line
\newcommand{\cvlabel}[2]{\EN{\textit{#1} }\ZH{\textit{#2}}}
